\documentclass[conference]{IEEEtran}
\IEEEoverridecommandlockouts

\usepackage{cite}
\usepackage{amsmath,amssymb,amsfonts}
\usepackage{graphicx}
\usepackage{textcomp}
\usepackage{xcolor}
\usepackage{booktabs}
\usepackage{url}
\usepackage{hyperref}
\hypersetup{colorlinks=true,linkcolor=black,citecolor=black,urlcolor=blue}

\def\BibTeX{{\rm B\kern-.05em{\sc i\kern-.025em b}\kern-.08em
    T\kern-.1667em\lower.7ex\hbox{E}\kern-.125emX}}

\begin{document}

\title{Lightweight Feature Selection for IDK Cascade Skip Decisions}

\author{
\IEEEauthorblockN{Kobe Young}
\IEEEauthorblockA{\textit{Department of Computer Science} \\
\textit{University of Houston}\\
Houston, TX, USA \\
kobe.tyler.young@gmail.com}
}

\maketitle

\begin{abstract}
IDK cascades chain progressively heavier classifiers so easy inputs exit early. Two recent improvements skip the middle classifier when the first stage is unlikely to be refined: Nguyen et al.\ (RTSS 2024) used a static confidence threshold; Katikaneni et al.\ (RTSS 2025) replaced it with a Random Forest on three softmax features. We ask whether a much simpler predictor suffices. Under Katikaneni et al.'s cross-dataset protocol (train on ImageNet-V2 Matched-Frequency + TopImages, test on Threshold-0.7) and their prediction target (will the middle classifier abstain?), a two-feature logistic regression on \textit{confidence} and \textit{margin} matches the Random Forest on the first-stage-IDK test population: XXlrlat vs.\ XXrflat~ms/image at XXlracc\% vs.\ XXrfacc\% top-1 accuracy (paired bootstrap 95\% CIs of the per-image differences include zero), at a measured XXratio$\times$ lower median per-decision evaluation cost (XXlrcost vs.\ XXrfcost at batch size 1 on a Jetson Orin Nano CPU). Unlike the tree ensemble, the deployed logistic is branch-free and data-independent, so its WCET is immediate to bound. For real-time deployments where the skip predictor's WCET itself enters the schedulability budget, a two-coefficient interpretable rule buys essentially the same cascade performance at a small fraction of the predictor overhead.
\end{abstract}

\begin{IEEEkeywords}
real-time machine learning, IDK cascades, anytime classification, edge inference, schedulability
\end{IEEEkeywords}

% =====================================================================
\section{Introduction}
\label{sec:intro}

Modern real-time systems increasingly rely on deep neural networks for perception, control, and classification at the edge. The dominant tension in this setting is well known: more accurate networks have larger and more variable inference times, which makes them difficult to schedule under hard deadlines. \textit{IDK cascades}, introduced by Wang et al.\ \cite{wang2017idk} and later formalized for real-time schedulability analysis by Baruah et al.\ \cite{baruah2023optimal}, address this tension by chaining a sequence of classifiers $A \to B \to C$ of increasing accuracy and cost. At each stage, the classifier either commits to a prediction or returns ``IDK'' (I-Don't-Know), in which case the next, heavier classifier is invoked. Easy inputs exit at $A$; only difficult inputs incur the full $A+B+C$ cost.

A recurring observation in the IDK-cascade literature is that the middle classifier $B$ is often \textit{wasted execution}: when $A$ already returns IDK with very low confidence, $B$ is unlikely to do better, and the cascade should jump directly to the deterministic final classifier $C$. Nguyen, Cheng, and Carroll \cite{nguyen2024probabilistic} formalized this observation with a probabilistic profiling phase and a single static threshold (skip $B$ when $A$'s top-1 confidence falls below 0.3), reporting 11.5--17.0\% average latency reduction across three ImageNet-V2 subsets with no loss of accuracy. Katikaneni, Cheng, and Carroll \cite{katikaneni2025randomforest} argued that a single static threshold does not generalize across datasets, and replaced it with a Random Forest trained on three features of $A$'s softmax distribution: \textit{confidence}, \textit{entropy}, and \textit{margin}. They showed that the learned skip predictor outperforms the static threshold on three ImageNet-V2 subsets and on a Jetson Nano embedded target.

The Random Forest, however, introduces its own per-image inference cost. The 2025 paper itself flags as future work the question of whether a simpler tabular model would suffice. In this work-in-progress, we take that question literally and ask:

\begin{quote}
\textit{Is a two-feature logistic skip rule sufficient to recover most of the latency benefit reported in \cite{katikaneni2025randomforest}, while remaining interpretable and trivially cheap to evaluate at inference time?}
\end{quote}

\textbf{Contributions.} (i) We reproduce the dynamic-threshold cascade of \cite{nguyen2024probabilistic} on a freely available ImageNet-V2 subset using off-the-shelf pre-trained ResNet-18, ResNet-34, and ResNet-152 weights. (ii) We implement a logistic-regression skip predictor over $\{$\textit{confidence}, \textit{margin}$\}$ and compare it to both the static threshold and a re-implementation of the Random Forest baseline with the same hyperparameters, under the cross-dataset protocol of \cite{katikaneni2025randomforest} (Sec.~\ref{sec:setup}), with bootstrap interval estimates and a nested-feature ablation. (iii) We measure the per-decision overhead of the skip predictor itself on an embedded ARM target and argue that, beyond its measured cost, the logistic's branch-free evaluation is the more analyzable term to carry into a schedulability bound.

% =====================================================================
\section{Background and Related Work}
\label{sec:related}

\subsection{IDK Cascades}
A formal IDK cascade \cite{baruah2023optimal} is a sequence of classifiers $C_1, C_2, \dots, C_n$ ordered by non-decreasing average execution time $\bar{e}_i$ and non-decreasing accuracy $a_i$. Each $C_i$ for $i < n$ is allowed to abstain by emitting an ``IDK'' token, which we denote $\bot$; the final $C_n$ must always commit. Let $p_i \in [0,1]$ denote the probability that $C_i$ emits $\bot$ on a randomly drawn input. The cascade's expected per-input response time is
\[
\mathbb{E}[R] \;=\; \sum_{i=1}^{n} \bar{e}_i \prod_{j<i} p_j ,
\]
and the schedulability question is whether the \emph{worst-case} response time, $R_{\text{wc}} = \sum_{i=1}^{n} e_i^{\text{wc}}$, satisfies $R_{\text{wc}} \le D$ for a deadline $D$. Baruah et al.\ \cite{baruah2023optimal} give an optimal-synthesis algorithm for cascade ordering when both $\bar{e}_i$ and $p_i$ are known. Abdelzaher et al.\ \cite{abdelzaher2023scheduling} extend the framework to scheduling cascades with arbitrary inter-task dependencies. Both works treat the cascade as a static structure: the decision of whether to invoke $C_{i+1}$ after an IDK at $C_i$ is fixed at design time.

\subsection{Static Threshold Skipping (Nguyen et al., 2024)}
Nguyen, Cheng, and Carroll \cite{nguyen2024probabilistic} observed that on ImageNet-V2 TopImages, when ResNet-18 returns a top-1 softmax probability below 0.3, ResNet-34 \emph{itself outputs IDK} in more than 92\% of cases (their Fig.~1). Since time is only saved by skipping $B$ when $B$ would output IDK, their dynamic cascade skips ResNet-34 whenever ResNet-18's confidence falls below 0.3, jumping directly to ResNet-152. Across three ImageNet-V2 subsets they report 11.5--17.0\% average latency reduction with no measurable accuracy degradation; on Matched-Frequency, the subset we use here, the reduction is 13.8\% (22.74 to 19.61~ms/image on a GTX~1070~Ti).

\subsection{Random Forest Skipping (Katikaneni et al., 2025)}
Katikaneni, Cheng, and Carroll \cite{katikaneni2025randomforest} replaced the single threshold with a Random Forest classifier (\texttt{n\_estimators=50}, \texttt{max\_depth=4}, \texttt{class\_weight='balanced'}) trained on three features of the ResNet-18 softmax output:
\begin{itemize}
    \item \textbf{Confidence}: $\max_i s_i$, the top-1 softmax probability.
    \item \textbf{Entropy}: $-\sum_i s_i \log s_i$, the uncertainty of the full distribution.
    \item \textbf{Margin}: $s_{(1)} - s_{(2)}$, the gap between top-1 and top-2 probabilities.
\end{itemize}
The label is whether ResNet-34 will itself output IDK. They train on the samples where $A$ abstained, drawn from the Matched-Frequency and TopImages subsets, and test on the Threshold-0.7 subset. On a Jetson Nano 4~GB embedded target (CPU mode), Random Forest skipping beat both the no-skip and static-threshold cascades on all three ImageNet-V2 subsets; on a desktop-class CPU the static threshold retained an edge on one subset. They suggest evaluating additional tabular models (e.g., XGBoost) as future work.

\subsection{Adjacent Real-Time ML Work}
Phonotomizer \cite{yantosca2025phonotomizer} ablates the stages of an online phonetic-segmentation pipeline to meet a deadline; the shared theme with IDK cascades is that the controller (skip predictor here, ablation manager there) must itself be cheap to evaluate, since its WCET is charged against the schedulability bound of \cite{baruah2023optimal, abdelzaher2023scheduling}.

% =====================================================================
\section{Approach}
\label{sec:approach}

\subsection{Cascade Architecture}
We use the canonical three-stage cascade from \cite{nguyen2024probabilistic, katikaneni2025randomforest}:
\[
A = \text{ResNet-18}\;\to\;B = \text{ResNet-34}\;\to\;C = \text{ResNet-152},
\]
all loaded with default \texttt{torchvision} pretrained weights (\texttt{IMAGENET1K\_V1}). A classifier $C_i$ emits $\bot$ (IDK) when its top-1 softmax probability is below an accuracy threshold $\tau_{\text{acc}} = 0.65$, matching the Nguyen et al.\ setup. The final classifier $C$ never abstains.

\subsection{Dynamic Skip Decision}
The control variable is $\mathrm{skip}(\mathbf{x}) \in \{0, 1\}$, evaluated only when $A$ returns $\bot$. Cascade execution: run $A$ and commit if $\max\mathbf{s}_A \ge \tau_{\text{acc}}$; otherwise consult $\mathrm{skip}$. If $\mathrm{skip}=1$, jump directly to $C$ and commit. If $\mathrm{skip}=0$, run $B$ and commit if $\max\mathbf{s}_B \ge \tau_{\text{acc}}$, else fall through to $C$.

\subsection{Skip-Decision Variants}
We compare four implementations of $\mathrm{skip}(\cdot)$:
\begin{enumerate}
    \item \textbf{Baseline}: $\mathrm{skip}(\cdot) \equiv 0$. Always run $B$ after $A$ abstains.
    \item \textbf{Static threshold} \cite{nguyen2024probabilistic}: $\mathrm{skip}(\mathbf{x}) = 1$ iff $\mathrm{conf}(A) < 0.3$.
    \item \textbf{Random Forest} (re-implementation of \cite{katikaneni2025randomforest} with the same hyperparameters, under our single-dataset protocol): a 50-tree, depth-4 RF classifier on three features $\{\text{conf}, \text{entropy}, \text{margin}\}$, trained with \texttt{class\_weight='balanced'} and \texttt{min\_samples\_leaf=40}.
    \item \textbf{Logistic (ours)}: a logistic regression on two features $\{\text{conf}, \text{margin}\}$ only.
\end{enumerate}

\subsection{Logistic Skip Rule}
Let $\mathbf{s}_A = (s_1, \dots, s_K) \in \Delta^{K-1}$ denote $A$'s softmax output and $s_{(1)} \ge s_{(2)} \ge \dots$ its order statistics. Define
\begin{align*}
\mathrm{conf}(A)   &= s_{(1)}, &
\mathrm{margin}(A) &= s_{(1)} - s_{(2)} .
\end{align*}
The skip predictor is
\[
P(\mathrm{skip}=1 \mid \mathbf{x}) = \sigma\!\left(w_0 + w_1\,\mathrm{conf}(A) + w_2\,\mathrm{margin}(A)\right),
\]
with $\sigma(z) = (1 + e^{-z})^{-1}$. The training label is $y = 1$ when $B$ itself would abstain, i.e.\ $\max \mathbf{s}_B < \tau_{\text{acc}}$. This is the predictor target of Katikaneni et al.\ \cite{katikaneni2025randomforest} and the waste condition identified by Nguyen et al.\ \cite{nguyen2024probabilistic}: time is only saved by skipping $B$ when $B$ would output IDK, whereas if $B$ would commit, running it is the cheapest path to a committed answer ($t_B < t_C$). Both learned predictors (RF and logistic) are trained on this label. We fit on the subset of images where $A$ returned $\bot$ (the only inputs where $\mathrm{skip}$ is consulted), drawn from the two training subsets under the cross-dataset protocol of Section~\ref{sec:setup}, and apply \texttt{class\_weight='balanced'} to handle the asymmetric prior ($y=1$ for XXtrainprior\% of $A$-IDK training images). At inference, the cascade skips $B$ when $P(\mathrm{skip}=1\mid \mathbf{x}) > 0.5$. Evaluating the predictor itself reduces to two multiplications, two additions, and one sigmoid; Section~\ref{sec:results} reports the measured per-decision cost of this rule against the 50-tree Random Forest.

% =====================================================================
\section{Experimental Setup}
\label{sec:setup}

\subsection{Datasets and Cross-Dataset Protocol}
We use all three ImageNet-V2 subsets \cite{recht2019imagenetv2} (10\,000 images each, 1\,000 classes): \textit{Matched-Frequency}, \textit{TopImages}, and \textit{Threshold-0.7}, obtained as the original per-variant release files of Recht et al.\ via the Hugging Face Hub, so the variant identity of each split is exact by construction. The standalone top-1 accuracies of our three classifiers on each subset are consistent with canonical ImageNet-V2 numbers. Following the cross-dataset protocol of Katikaneni et al.\ \cite{katikaneni2025randomforest}, both learned predictors are trained on the first-stage-IDK samples of Matched-Frequency and TopImages, and all evaluation is performed on Threshold-0.7. Because training and test data come from disjoint datasets, no within-dataset split is needed and the full $A$-IDK population of the test subset is used for evaluation; this also directly tests the transfer property that motivated \cite{katikaneni2025randomforest} to replace the static threshold.

\subsection{Hardware}
All inference is run in CPU mode (\texttt{torch.device("cpu")}), matching the GPU-disabled methodology of Katikaneni et al.\ \cite{katikaneni2025randomforest}. CPU is the operative regime: a single ResNet-152 GPU run is already cheap, so the cascade's value vanishes there. Our primary platform is an NVIDIA Jetson Orin Nano (ARM Cortex-A78AE, CPU mode), the successor of the Jetson Nano used by \cite{katikaneni2025randomforest}; softmax outputs are platform-independent and were cached on an x86\_64 host, while all reported per-image stage timings and per-decision predictor costs are measured on the board.

\subsection{Implementation}
Models are loaded from \texttt{torchvision.models} with default pretrained weights. Per-image timings are measured with \texttt{time.perf\_counter()} around the inference call only (preprocessing is excluded). Each measurement is the mean of five trials per image after a 50-image warm-up.

\subsection{Metrics and Evaluation Populations}
\begin{itemize}
    \item \textbf{Average latency} (ms/image): mean wall-clock time per image over the evaluation population stated with each table.
    \item \textbf{Top-1 accuracy} on the cascade's final committed prediction, over the same population.
    \item \textbf{Skip rate}: fraction of images in the population on which the cascade skipped $B$.
    \item For the learned predictors: precision, recall, F1 of the skip-decision binary classifier, and measured per-decision prediction cost.
\end{itemize}
Two populations appear in Section~\ref{sec:results}. The threshold sweep uses the \emph{full 10\,000-image Matched-Frequency set} (a reproduction of \cite{nguyen2024probabilistic}). The learned-predictor comparisons use the \emph{$A$-IDK population of Threshold-0.7}: the skip rule is consulted nowhere else, and no Threshold-0.7 image appears in training. Numbers from the two populations are not directly comparable (the $A$-IDK population contains only hard inputs, so its latencies are higher and accuracies lower than full-set values). All point estimates on the test population are reported with bootstrap 95\% confidence intervals (10\,000 percentile-bootstrap resamples of the test population); comparisons between the Random Forest and the logistic use \emph{paired} bootstrap CIs of the per-image differences.

% =====================================================================
\section{Preliminary Results}
\label{sec:results}

All numbers below are measured on the ImageNet-V2 \emph{Matched-Frequency} subset on x86 CPU, with five trials per image after a 50-image warm-up. Table~\ref{tab:threshold-sweep} is computed on the full 10\,000-image set; Tables~\ref{tab:variants}--\ref{tab:overhead} are computed on the held-out $A$-IDK test split defined in Section~\ref{sec:setup}.

\subsection{Threshold Sweep Reproduction}
Table~\ref{tab:threshold-sweep} reproduces the threshold-sweep protocol of Nguyen et al.\ \cite{nguyen2024probabilistic}: cascade latency as a function of the ResNet-18 confidence threshold $\tau$ used for the skip-to-$C$ rule, on the full 10\,000-image set. The minimum sits at $\tau=0.4$ (73.68 ms), with the 0.3 setting used in the original paper a close second (74.06 ms); both beat the no-skip baseline (77.49 ms), a 4.4--4.9\% reduction. The relative gain is smaller than the 13.8\% Nguyen et al.\ report on Matched-Frequency, which we attribute to the hardware difference (their timings are GTX~1070~Ti GPU; ours are CPU, with a different stage-cost ratio). That our optimum lands at 0.4 rather than their 0.3 is consistent with Katikaneni et al.'s argument \cite{katikaneni2025randomforest} that a single static threshold does not transfer across datasets. Accuracy is flat to slightly increasing across the sweep (0.651 to 0.656): below the threshold region $B$ nearly always abstains \cite{nguyen2024probabilistic}, so replacing $B$ with the more accurate $C$ never hurts the committed prediction.

\begin{table}[h]
\centering
\caption{Threshold Sweep on the Full 10\,000-Image Set (Protocol of \cite{nguyen2024probabilistic})}
\label{tab:threshold-sweep}
\begin{tabular}{cccc}
\toprule
$\tau$ & Latency (ms) & Accuracy & Skip rate \\
\midrule
0.0 & 77.49 & 0.651 & 0.000 \\
0.1 & 76.96 & 0.651 & 0.018 \\
0.2 & 75.36 & 0.651 & 0.089 \\
0.3 & 74.06 & 0.652 & 0.181 \\
0.4 & 73.68 & 0.653 & 0.272 \\
0.5 & 74.07 & 0.653 & 0.363 \\
0.6 & 75.40 & 0.655 & 0.448 \\
0.7 & 76.20 & 0.656 & 0.486 \\
0.8 & 76.20 & 0.656 & 0.486 \\
0.9 & 76.20 & 0.656 & 0.486 \\
\bottomrule
\end{tabular}

\end{table}

\subsection{Skip-Predictor Comparison}
Table~\ref{tab:variants} reports our headline comparison on the Threshold-0.7 $A$-IDK population (XXntest images), with the predictors trained on Matched-Frequency + TopImages. The Random Forest re-implementation recovers Katikaneni et al.'s qualitative result under cross-dataset transfer: a learned predictor beats the static threshold (XXrflat vs.\ XXstaticlat ms). The two-feature logistic rule (\textbf{ours}) matches the Random Forest on cascade behavior: the paired bootstrap 95\% CIs of the per-image RF$-$logistic differences are $[-2.62, -0.28]$~ms for latency and $[-0.002, +0.002]$ for accuracy
\hspace{-0.3em}, so neither predictor dominates the other on either metric. Predictor quality is compared in Table~\ref{tab:skip-f1} with bootstrap CIs on F1.

Point estimates at this scale are fragile: in a within-subset control experiment (10 random 70/30 splits of the Matched-Frequency $A$-IDK samples), skip-F1 varied by $\pm 0.02$ across splits for the Random Forest (0.677~$\pm$~0.017) against 0.704~$\pm$~0.008 for the logistic --- the split-to-split spread exceeds the gap between the two predictors, which is why we report interval estimates throughout.

Table~\ref{tab:ablation} gives the feature-selection evidence behind the two-feature choice: nested logistic models on $\{$conf$\}$, $\{$conf, margin$\}$, and $\{$conf, margin, entropy$\}$. XXablationsummary

\begin{table}[h]
\centering
\caption{Cascade Variants on the Threshold-0.7 $A$-IDK Population (Cross-Dataset Protocol; Bootstrap 95\% CIs)}
\label{tab:variants}
\begin{tabular}{lccc}
\toprule
Variant & Latency (ms) & Accuracy & Skip rate \\
\midrule
Baseline (no skip) & 374.86\,{\scriptsize$[369.91,\,379.82]$} & 0.549\,{\scriptsize$[0.533,\,0.564]$} & 0.000 \\
Static $\tau=0.3$ & 363.65\,{\scriptsize$[359.30,\,368.01]$} & 0.552\,{\scriptsize$[0.537,\,0.567]$} & 0.315 \\
Random Forest & 363.83\,{\scriptsize$[360.03,\,367.63]$} & 0.553\,{\scriptsize$[0.538,\,0.568]$} & 0.512 \\
\textbf{Logistic (ours)} & 365.28\,{\scriptsize$[361.48,\,369.05]$} & 0.553\,{\scriptsize$[0.538,\,0.568]$} & 0.493 \\
\bottomrule
\end{tabular}

\end{table}

\begin{table}[h]
\centering
\caption{Skip-Predictor Quality on the Threshold-0.7 $A$-IDK Population}
\label{tab:skip-f1}
\begin{tabular}{lccc}
\toprule
Predictor & Precision & Recall & F1 \\
\midrule
RF \{conf, entropy, margin\} & 0.778 & 0.607 & 0.682\,{\scriptsize$[0.667,\,0.697]$} \\
Logistic \{conf\} & 0.778 & 0.598 & 0.676\,{\scriptsize$[0.661,\,0.692]$} \\
\textbf{Logistic \{conf, margin\} (ours)} & 0.772 & 0.580 & 0.662\,{\scriptsize$[0.647,\,0.678]$} \\
Logistic \{conf, margin, entropy\} & 0.772 & 0.579 & 0.662\,{\scriptsize$[0.647,\,0.678]$} \\
\bottomrule
\end{tabular}

\end{table}

\begin{table}[h]
\centering
\caption{Logistic Skip-F1 Under Nested Feature Sets (Cross-Dataset Protocol)}
\label{tab:ablation}
\input{table_ablation}
\end{table}

\subsection{Measured Predictor Overhead}
Table~\ref{tab:overhead} reports the per-decision cost of each learned predictor, measured at batch size 1 over 1\,000 test samples after warm-up (\texttt{time.perf\_counter}, on the Jetson Orin Nano): the scikit-learn Random Forest via \texttt{predict()}, the logistic in its deployed form (a NumPy dot product plus sigmoid). The median gap is XXratio$\times$ (XXrfcost vs.\ XXlrcost). Because worst-case behavior is what enters a schedulability bound, we also report the empirical p99 and maximum. This measures the deployed scikit-learn implementation --- a compiled tree ensemble would narrow the gap --- but the logistic's two-multiply evaluation is a floor no 50-tree ensemble reaches. Beyond magnitude, the two predictors differ in \emph{analyzability}: the deployed logistic is branch-free and data-independent (exactly two multiplications, two additions, and one sigmoid on every input), so a WCET bound is immediate, whereas a tree ensemble's evaluation path is input-dependent and its interpreted implementation resists tight static bounding. The cascade latencies in Table~\ref{tab:variants} \emph{exclude} predictor cost; on the $A$-IDK population, where the predictor is consulted on every image, charging each variant its median per-decision cost would add XXrfcost to the Random Forest row but only XXlrcostms to the logistic row.

\begin{table}[h]
\centering
\caption{Measured Per-Decision Predictor Cost (Batch Size 1, Jetson Orin Nano CPU, 1\,000 Test Samples)}
\label{tab:overhead}
\begin{tabular}{lccc}
\toprule
Predictor & Mean ($\mu$s) & Median ($\mu$s) & Relative \\
\midrule
Random Forest & 3278.1 & 2795.6 & $84\times$ \\
\textbf{Logistic (ours)} & 35.2 & 33.2 & $1\times$ \\
\bottomrule
\end{tabular}

\end{table}

\paragraph{Learned coefficients.} The logistic predictor's two-feature linear discriminant is, in the trained model,
\[ \mathrm{logit}\,P(\mathrm{skip}\mid \mathbf{x}) = 1.787 -5.111\,\mathrm{conf}(A) +1.086\,\mathrm{margin}(A) . \]

The dominant term matches intuition: higher confidence sharply reduces skip probability (a confident $A$ suggests an input easy enough for $B$ to commit on). The smaller positive margin coefficient applies a correction conditional on confidence: at fixed top-1 confidence, a larger top-1/top-2 gap slightly increases the predicted probability that $B$ abstains. Confidence dominates --- consistent with confidence being the single feature used by the static threshold rule.

% =====================================================================
\section{Discussion and Future Work}
\label{sec:discussion}

\textbf{Why two features suffice.} The skip decision is approximately a function of the top-2 ordering of the softmax. On a $K{=}1000$-class output, the bulk of entropy lives in the long tail of small probabilities, but whether $B$ will abstain is dominated by top-1/top-2 ambiguity --- already captured by confidence and margin. The nested-model ablation (Table~\ref{tab:ablation}) makes this concrete: XXablationdiscussion

\textbf{Real-time implications.} The skip predictor's own WCET enters the schedulability budget. Our measured per-decision cost gap is XXratio$\times$ at the median (Table~\ref{tab:overhead}), and the logistic's branch-free, data-independent evaluation admits an immediate WCET bound where the tree ensemble's input-dependent path does not; in Baruah's framing \cite{baruah2023optimal}, where each stage's WCET is charged against the response-time bound, the near-zero-cost predictor is preferable whenever cascade accuracy and latency are comparable --- which Table~\ref{tab:variants} shows they are.

\textbf{Future work.} Confidence and margin are well-defined for any softmax classifier, so the logistic rule is dataset-agnostic in form (e.g., audio via Phonotomizer \cite{yantosca2025phonotomizer}, tabular medical data). Natural follow-ons: fold the predictor's WCET into Baruah's response-time bound, and generalize the binary skip to a multinomial multi-stage rule for cascades with more than three classifiers.

% =====================================================================
\section{Conclusion}
\label{sec:conclusion}

We tested the future-work question raised by Katikaneni et al.\ \cite{katikaneni2025randomforest}: does a much simpler skip predictor recover most of the latency benefit? Under their cross-dataset protocol (train on Matched-Frequency + TopImages, test on the Threshold-0.7 $A$-IDK population), a two-feature logistic regression on confidence and margin matches the Random Forest on cascade latency (XXlrlat vs.\ XXrflat~ms) and accuracy (XXlracc\% vs.\ XXrfacc\%) --- the paired bootstrap CIs of the per-image differences include zero --- at a measured XXratio$\times$ lower median per-decision cost on a Jetson Orin Nano, with a branch-free evaluation whose WCET is immediate to bound. For real-time deployments where the skip predictor's WCET is a first-class scheduling concern, the two-coefficient rule buys the same cascade performance at a small fraction of the predictor overhead. Folding that bound into an end-to-end schedulability analysis of the full cascade is the immediate next step.

% =====================================================================
\bibliographystyle{IEEEtran}
\bibliography{references}

\end{document}
